%=============================================================================
% KAPPA_{U(1)} FORMULA v2: DEFINITIVE ANALYSIS
% Exhaustive search of all DFD source files, March 2026
%=============================================================================
\documentclass[12pt,a4paper]{article}
\usepackage[margin=2.5cm]{geometry}
\usepackage{amsmath,amssymb,amsthm}
\usepackage{booktabs,array,longtable}
\usepackage{hyperref}
\usepackage{xcolor}

\newtheorem{theorem}{Theorem}
\newtheorem{proposition}{Proposition}
\newtheorem{finding}{Finding}
\theoremstyle{definition}
\newtheorem{definition}{Definition}
\newtheorem{gap}{Gap}

\title{The Complete Formula for $\kappa_{U(1)}$ and $\alpha^{-1}$:\\
Definitive Extraction from the DFD Corpus (v2)}

\author{Extracted from G.\ Alcock, \textit{Density Field Dynamics}\\
All source files exhaustively searched}
\date{March 22, 2026 --- v2}

\begin{document}
\maketitle

\begin{abstract}
After exhaustive line-by-line reading of every \texttt{.tex} file in the
DFD v108 source tree, the standalone ``Ab Initio'' PDF, and all
supplementary notes, this document presents the \textbf{definitive answer}
to the question: \emph{What is the explicit closed-form algebraic expression
$\kappa_{U(1)} = G(k_{\max}, \beta_{U(1)}, d, w, \varepsilon_{\mathrm{adj}}, \ldots)$
that someone can plug into a calculator to get
$\kappa_{U(1)} = 7.26999\ldots \to \alpha^{-1} = 6\pi\kappa_{U(1)} = 137.03599985$?}

The answer has three layers of increasing precision:
\begin{enumerate}
\item[\textbf{Layer 1}] (Exact): $\alpha^{-1} = 6\pi\,\kappa_{U(1)}$ ---
  this is a rigorous algebraic identity from electroweak mixing $+$
  the Frame Stiffness Theorem.
\item[\textbf{Layer 2}] (85~ppm approximate):
  $\alpha^{-1} \approx \frac{35\pi}{3}\,\beta_{U(1)}\,
  \frac{d-1}{d}\,\frac{d^2}{d^2-1}$
  where $\beta_{U(1)} = 3.796945\ldots$ is a computable finite sum.
\item[\textbf{Layer 3}] (Sub-ppm): The paper claims this comes from
  numerical evaluation of the $a_4$ Seeley--DeWitt coefficient on the
  Toeplitz-truncated $\mathbb{C}P^2 \times S^3$.  \textbf{No closed-form
  algebraic expression for this coefficient appears anywhere in the
  DFD corpus.}  The spectral action computation is referenced but
  never written out.
\end{enumerate}
\end{abstract}

\tableofcontents
\newpage

%=============================================================================
\section{Layer 1: The Exact Electroweak Formula}
\label{sec:layer1}
%=============================================================================

\begin{theorem}[Master formula --- exact]
\label{thm:master}
\begin{equation}
\boxed{\alpha^{-1} = 6\pi\,\kappa_{U(1)}}
\end{equation}
\end{theorem}

\begin{proof}
From the DFD gauge-emergence stiffness functional, the gauge couplings
are extracted via Wilson normalization conventions (Ab Initio paper, Eqs.~12--13):
\begin{equation}
g_1^2 = \frac{1}{\kappa_{U(1)}}\,,\qquad
g_2^2 = \frac{4}{\kappa_{SU(2)}}\,.
\end{equation}

Electroweak mixing gives the electromagnetic coupling:
\begin{equation}
\frac{1}{e^2} = \frac{1}{g_1^2} + \frac{1}{g_2^2}
= \kappa_{U(1)} + \frac{\kappa_{SU(2)}}{4}\,.
\end{equation}

The Frame Stiffness Theorem (derived from gauge-emergence geometry on
$\mathbb{C}P^2 \times S^3$ with partition $(3,2,1)$) gives:
\begin{equation}
\frac{\kappa_{U(1)}}{\kappa_{SU(2)}} = \frac{n_1}{n_2} = \frac{1}{2}\,,
\qquad\text{so}\qquad \kappa_{SU(2)} = 2\,\kappa_{U(1)}\,.
\end{equation}

Substituting:
\begin{equation}
\frac{1}{e^2} = \kappa_{U(1)} + \frac{2\,\kappa_{U(1)}}{4}
= \frac{3}{2}\,\kappa_{U(1)}\,.
\end{equation}

With $\alpha = e^2/(4\pi)$:
\begin{equation}
\alpha^{-1} = \frac{4\pi}{e^2} = 4\pi \cdot \frac{3}{2}\,\kappa_{U(1)}
= 6\pi\,\kappa_{U(1)}\,. \qed
\end{equation}
\end{proof}

\paragraph{Numerical requirement.}
The experimental value $\alpha^{-1}_{\mathrm{exp}} = 137.035999084(21)$ requires:
\begin{equation}
\kappa_{U(1)} = \frac{137.035999084}{6\pi} = 7.269986\ldots
\end{equation}

The DFD prediction $\alpha^{-1} = 137.03599985$ corresponds to
$\kappa_{U(1)} = 7.269986\ldots$ (residual $-0.006$~ppm).

%=============================================================================
\section{Layer 2: The Best Approximate Closed-Form Expression}
\label{sec:layer2}
%=============================================================================

\subsection{The computable input: $\beta_{U(1)}$}

The bare lattice coupling is a \emph{finite, computable sum}:

\begin{definition}[CS weighted average]
\begin{equation}
\boxed{
\beta_{U(1)} = \frac{\displaystyle\sum_{k=0}^{k_{\max}-1}(k+2)\,w(k)}
     {\displaystyle\sum_{k=0}^{k_{\max}-1}w(k)}\,,
\qquad
w(k) = \frac{2}{k+2}\sin^2\!\frac{\pi}{k+2}\,,
\qquad k_{\max} = 60\,.
}
\end{equation}
\end{definition}

The weights $w(k) = |Z_{SU(2)_k}(S^3)|^2$ are the squared $SU(2)_k$
Chern--Simons partition functions on $S^3$.  The shifted level $k+2$
arises from the dual Coxeter number $h^\vee = 2$ for $SU(2)$.

\paragraph{Numerical evaluation.}
\begin{equation}
\beta_{U(1)}\big|_{k_{\max}=60} = 3.796945275657\ldots \approx 3.80\,.
\end{equation}

\subsection{The approximate formula}

\begin{proposition}[Approximate closed form --- 85~ppm]
\label{prop:approx}
\begin{equation}
\boxed{
\alpha^{-1} \approx \frac{35\pi}{3}\;\beta_{U(1)}\;\frac{d-1}{d}
\;\frac{d^2}{d^2-1}
}
\end{equation}
where $d = k_{\max} + 4 = 64$.  This evaluates to $137.024$, accurate to 85~ppm.
\end{proposition}

\paragraph{Factorization.}
The prefactor $35\pi/3$ decomposes as:
\begin{equation}
\frac{35\pi}{3} = 4\pi \times \underbrace{\frac{5}{2}}_{\text{EW mixing}}
\times \underbrace{\frac{7}{6}}_{\text{spectral triple}}\,,
\end{equation}
where:
\begin{itemize}
\item $\frac{5}{2} = 1 + \frac{R_W}{4}$ with Wilson ratio $R_W = 6$
  encodes the electroweak mixing
  $1/e^2 = \beta + \beta \cdot R_W/4 = (5/2)\,\beta$ at tree level.
\item $\frac{7}{6} = \frac{N_{\mathrm{species}}}{N_{\mathrm{gen}} \cdot n_2}
  = \frac{7}{3 \times 2}$ encodes the microsector spectral triple correction.
\end{itemize}

\paragraph{Numerical verification.}
\begin{align}
\frac{35\pi}{3} \times 3.796945 \times \frac{63}{64} \times \frac{4096}{4095}
&= 36.6519\ldots \times 3.796945 \times 0.984375 \times 1.000244 \notag\\
&= 137.024\ldots \notag\\
&= 137.036 - 0.012 \qquad (\text{residual: } -85\text{ ppm}).
\end{align}

\subsection{Why the approximate formula is not exact}

The ratio $\kappa_{U(1)}/\beta_{U(1)} = 7.26999/3.79695 = 1.91469\ldots$
and the tree-level prediction would be $35/18 = 1.94444\ldots$.
These differ by $1.6\%$, reflecting the non-perturbative renormalization
from bare coupling to renormalized stiffness.  The tree-level formula
cannot capture this; the full spectral action or Monte Carlo is required.

%=============================================================================
\section{Layer 3: The Sub-ppm Route (What the Paper Claims)}
\label{sec:layer3}
%=============================================================================

\subsection{The claim}

Section~8C of the Unified Theory (v108) states that $\alpha^{-1} = 137.03599985$
is obtained from the spectral action on the Toeplitz-truncated internal geometry
$\mathbb{C}P^2 \times S^3$, using the following eight locked inputs:

\begin{table}[h]
\centering
\caption{Complete input table (Table~XXVIII of the Unified Theory).}
\label{tab:inputs}
\renewcommand{\arraystretch}{1.3}
\begin{tabular}{llll}
\toprule
\textbf{Component} & \textbf{Value} & \textbf{Source} & \textbf{Status} \\
\midrule
$K_{\mathbb{C}P^2} = \mathcal{O}(-3)$ & $-3$
  & Algebraic geometry & Rigorous \\
$L_{\det} = K^{-1}$ & $\mathcal{O}(3)$
  & Spin$^c$ structure & Rigorous \\
$d = k_{\max} + 4$ & $64$
  & $\dim H^0(\mathcal{O}(k_{\max}+3))$ & Rigorous \\
$(d-1)/d$ & $63/64$
  & Traceless projection & Derived \\
$N_{\mathrm{species}}$ & $7$
  & SM SU(2) doublet components & SM content \\
$\mathrm{Tr}(Y^2)$ & $10$
  & SM hypercharge sum & SM content \\
$g_F$ & $8$
  & Spectral triple ($J \times \gamma \times \mathbb{C}$) & Derived \\
$w = N_{\mathrm{species}}/(g_F \cdot \mathrm{Tr}(Y^2))$ & $7/80$
  & Hypercharge weighting & Derived \\
$\varepsilon_{\mathrm{adj}} = d^2/(d^2-1)$ & $4096/4095$
  & Regular-module BOOST & Forced \\
\bottomrule
\end{tabular}
\end{table}

The computation proceeds through:
\begin{enumerate}
\item \textbf{Operator:} Connection Laplacian $P = -g^{ij}\nabla_i\nabla_j$
  with parallel K\"ahler curvature (kills derivative ambiguities).
\item \textbf{Regulator:} Plateau cutoff with $f_{-2} = f_{-4} = \cdots = 0$
  (eliminates $a_6, a_8, \ldots$ as hidden tuning knobs).
\item \textbf{Finite-$k$ truncation:} Toeplitz quantization at level $m = k+3$
  gives $d = 64$ and spectral cutoff $\Lambda^3 = k(d{-}1)/d = 59.0625$.
\item \textbf{SM content:} Hypercharge weight $w = 7/80$ enters through the
  $a_4$ trace over the internal fermion Hilbert space.
\item \textbf{Microsector trace:} Regular-module $\mathcal{H}_F = M_d(\mathbb{C})$
  forces the BOOST factor $\varepsilon_{\mathrm{adj}} = 4096/4095$.
\end{enumerate}

\subsection{The gap: no explicit $a_4$ formula}

\begin{gap}[Missing spectral action formula]
\label{gap:a4}
\textbf{After exhaustive search of every source file}, no explicit formula
for the $a_4$ Seeley--DeWitt coefficient evaluated on the Toeplitz-truncated
$\mathbb{C}P^2 \times S^3$ geometry appears anywhere in the DFD corpus.

The standard heat-kernel formula
\begin{equation}
a_4(P) = \frac{1}{(4\pi)^{d/2}} \int_X \mathrm{tr}\!\left(
\frac{1}{180}(R_{\mu\nu\rho\sigma}R^{\mu\nu\rho\sigma}
- R_{\mu\nu}R^{\mu\nu} + \tfrac{5}{2}R^2)
+ \frac{1}{12}\Omega_{\mu\nu}\Omega^{\mu\nu}
+ \cdots \right) \mathrm{dvol}
\end{equation}
is never written out, nor evaluated on the specific geometry.
The paper says ``the gauge-kinetic extraction from $a_4$ is unambiguous''
but never shows the computation.
\end{gap}

\paragraph{What IS present.}
The paper presents the \emph{result} $\kappa_{U(1)} = 7.26999\ldots$
and $\alpha^{-1} = 137.03599985$ as outputs of this computation,
but the intermediate steps are not shown.

\subsection{The forced binary fork (what IS shown)}

The paper does show how the microsector trace choice splits the result:

\begin{center}
\renewcommand{\arraystretch}{1.3}
\begin{tabular}{lccc}
\toprule
Branch & Factor & $\alpha^{-1}$ & Residual (ppm) \\
\midrule
\textbf{A (regular-module)} & $4096/4095$ & $137.03599985$ & $-0.006$ \\
B (fermion-rep) & $4095/4096$ & $137.03014445$ & $+42.7$ \\
\midrule
Experimental & --- & $137.035999084$ & --- \\
\bottomrule
\end{tabular}
\end{center}

The two branches differ by a multiplicative factor of
$(4096/4095)/(4095/4096) = (4096)^2/(4095)^2 \approx 1 + 2/4095$.
This tells us the base stiffness (before the BOOST/DROP choice) is:
\begin{equation}
\kappa_{\mathrm{base}} = \frac{\alpha^{-1}_A}{6\pi \cdot \varepsilon_A}
= \frac{137.03599985}{6\pi \times (4096/4095)} = 7.26821\ldots
\end{equation}

The branch difference gives $G \approx 12 = 4 \times N_{\mathrm{gen}}$,
confirming the three-generation structure.

%=============================================================================
\section{The Complete Finite Sum for $\beta_{U(1)}$}
\label{sec:eigenvalues}
%=============================================================================

For completeness, here is the explicit computation. The 60 Chern--Simons
levels run $k = 0, 1, \ldots, 59$ with shifted eigenvalues
$\lambda_k = k+2$ from $2$ to $61$.

\subsection{The weight function}

\begin{equation}
w(k) = \frac{2}{k+2}\sin^2\!\frac{\pi}{k+2}\,,
\qquad k = 0, 1, \ldots, 59\,.
\end{equation}

This is the squared modulus of the $SU(2)_k$ CS partition function on $S^3$:
\begin{equation}
Z_{SU(2)_k}(S^3) = \sqrt{\frac{2}{k+2}}\,\sin\!\frac{\pi}{k+2}\,.
\end{equation}

\subsection{Numerical values (all 60 levels)}

{%
\small
\begin{longtable}{rr@{\quad}r@{\quad}r@{\quad}r}
\toprule
$k$ & $k{+}2$ & $w(k)$ & $(k{+}2)\cdot w(k)$ & Running $\beta$ \\
\midrule
\endhead
 0 &  2 & 1.000000000000 & 2.000000000000 & 2.000000 \\
 1 &  3 & 0.500000000000 & 1.500000000000 & 2.333333 \\
 2 &  4 & 0.250000000000 & 1.000000000000 & 2.571429 \\
 3 &  5 & 0.138196601125 & 0.690983005625 & 2.749175 \\
 4 &  6 & 0.083333333333 & 0.500000000000 & 2.886582 \\
 5 &  7 & 0.053787171163 & 0.376510198141 & 2.995824 \\
 6 &  8 & 0.036611652352 & 0.292893218813 & 3.084678 \\
 7 &  9 & 0.025995061876 & 0.233955556881 & 3.158325 \\
 8 & 10 & 0.019098300563 & 0.190983005625 & 3.220339 \\
 9 & 11 & 0.014431497015 & 0.158746467169 & 3.273261 \\
10 & 12 & 0.011164549685 & 0.133974596216 & 3.318947 \\
11 & 13 & 0.008811074950 & 0.114543974347 & 3.358780 \\
12 & 14 & 0.007073652293 & 0.099031132098 & 3.393815 \\
13 & 15 & 0.005763636157 & 0.086454542357 & 3.424867 \\
14 & 16 & 0.004757529218 & 0.076120467489 & 3.452577 \\
15 & 17 & 0.003972221800 & 0.067527770596 & 3.477456 \\
16 & 18 & 0.003350409956 & 0.060307379214 & 3.499916 \\
17 & 19 & 0.002851724121 & 0.054182758299 & 3.520293 \\
18 & 20 & 0.002447174185 & 0.048943483705 & 3.538864 \\
19 & 21 & 0.002115580677 & 0.044427194214 & 3.555857 \\
20 & 22 & 0.001841228472 & 0.040507026386 & 3.571467 \\
21 & 23 & 0.001612291854 & 0.037082712652 & 3.585854 \\
22 & 24 & 0.001419757238 & 0.034074173711 & 3.599158 \\
23 & 25 & 0.001256673555 & 0.031416838871 & 3.611495 \\
24 & 26 & 0.001117622407 & 0.029058182574 & 3.622967 \\
25 & 27 & 0.000998338127 & 0.026955129420 & 3.633663 \\
26 & 28 & 0.000895431708 & 0.025072087818 & 3.643658 \\
27 & 29 & 0.000806187734 & 0.023379444290 & 3.653019 \\
28 & 30 & 0.000728413309 & 0.021852399266 & 3.661805 \\
29 & 31 & 0.000660324476 & 0.020470058748 & 3.670066 \\
30 & 32 & 0.000600459987 & 0.019214719597 & 3.677849 \\
31 & 33 & 0.000547615234 & 0.018071302737 & 3.685194 \\
32 & 34 & 0.000500791186 & 0.017026900316 & 3.692136 \\
33 & 35 & 0.000459154611 & 0.016070411401 & 3.698709 \\
34 & 36 & 0.000422006861 & 0.015192246988 & 3.704940 \\
35 & 37 & 0.000388759180 & 0.014384089652 & 3.710856 \\
36 & 38 & 0.000358913068 & 0.013638696597 & 3.716479 \\
37 & 39 & 0.000332044548 & 0.012949737362 & 3.721832 \\
38 & 40 & 0.000307791485 & 0.012311659405 & 3.726933 \\
39 & 41 & 0.000285843322 & 0.011719576220 & 3.731800 \\
40 & 42 & 0.000265932709 & 0.011169173775 & 3.736448 \\
41 & 43 & 0.000247828649 & 0.010656631925 & 3.740891 \\
42 & 44 & 0.000231330866 & 0.010178558119 & 3.745144 \\
43 & 45 & 0.000216265139 & 0.009731931258 & 3.749217 \\
44 & 46 & 0.000202479434 & 0.009314053964 & 3.753123 \\
45 & 47 & 0.000189840678 & 0.008922511845 & 3.756870 \\
46 & 48 & 0.000178232055 & 0.008555138626 & 3.760470 \\
47 & 49 & 0.000167550738 & 0.008209986177 & 3.763929 \\
48 & 50 & 0.000157705974 & 0.007885298686 & 3.767257 \\
49 & 51 & 0.000148617457 & 0.007579490328 & 3.770460 \\
50 & 52 & 0.000140213960 & 0.007291125902 & 3.773546 \\
51 & 53 & 0.000132432151 & 0.007018903986 & 3.776521 \\
52 & 54 & 0.000125215597 & 0.006761642258 & 3.779390 \\
53 & 55 & 0.000118513903 & 0.006518264651 & 3.782160 \\
54 & 56 & 0.000112281966 & 0.006287790107 & 3.784835 \\
55 & 57 & 0.000106479345 & 0.006069322682 & 3.787420 \\
56 & 58 & 0.000101069704 & 0.005862042846 & 3.789919 \\
57 & 59 & 0.000096020335 & 0.005665199790 & 3.792337 \\
58 & 60 & 0.000091301744 & 0.005478104632 & 3.794678 \\
59 & 61 & 0.000086887285 & 0.005300124385 & 3.796945 \\
\midrule
\multicolumn{2}{l}{$\sum w(k)$} & $2.192417$ & & \\
\multicolumn{2}{l}{$\sum (k{+}2)\,w(k)$} & & $8.324487$ & \\
\multicolumn{2}{l}{$\beta_{U(1)}$} & & & $\mathbf{3.796945}$ \\
\bottomrule
\end{longtable}
}

\subsection{Python verification code}

\begin{verbatim}
import math
def w(k): return (2.0/(k+2)) * math.sin(math.pi/(k+2))**2
Z   = sum(w(k) for k in range(60))
num = sum((k+2)*w(k) for k in range(60))
beta = num / Z   # = 3.796945275657...
\end{verbatim}

%=============================================================================
\section{The Two Routes to $\kappa_{U(1)}$}
\label{sec:routes}
%=============================================================================

\subsection{Route A: Lattice Monte Carlo ($\sim 1\%$ precision)}

\begin{enumerate}
\item Input $(\beta_{U(1)}, \beta_{SU(2)}) = (3.80, 22.80)$ into
  the DFD lattice action (Ab Initio Eq.~19).
\item Run Metropolis Monte Carlo on $L^4$ lattice.
\item Measure $\kappa_{U(1)}$ via background-field method (Eq.~20):
  \begin{equation}
  \kappa = \frac{F''(0)}{V}\,,\qquad
  F''(0) = \langle S'' \rangle - \langle (S')^2 \rangle + \langle S' \rangle^2\,.
  \end{equation}
\item Extract $\alpha_W = 1/(6\pi\kappa_{U(1)})$.
\item Result: $\alpha^{-1} \approx 137 \pm 2$ (86 runs, $L=4$--16).
\end{enumerate}

\textbf{Key fact (Ab Initio Eq.~11):} ``The relationship between
[$\beta$ and $\kappa$] is non-perturbative and lattice-size dependent,
which is precisely why the simulation is needed.''

\subsection{Route B: Spectral action (sub-ppm, claimed)}

The Unified Theory, Section~8C, claims that the spectral action on the
Toeplitz-truncated $\mathbb{C}P^2 \times S^3$ computes $\kappa_{U(1)}$
analytically from the $a_4$ Seeley--DeWitt coefficient of the connection
Laplacian.

\textbf{What is NOT shown:} The explicit formula for $a_4$ on this geometry.

\textbf{What IS shown:} The eight input components (Table~\ref{tab:inputs}),
the plateau cutoff that kills $a_6$ and higher, the Toeplitz truncation that
gives $d = 64$, and the forced fork between BOOST ($4096/4095$) and DROP
($4095/4096$) microsector traces.

\textbf{The result} $\alpha^{-1} = 137.03599985$ is stated without showing
intermediate steps. The computation appears to have been carried out
numerically from the spectral action framework, not reduced to a single
algebraic expression.

%=============================================================================
\section{Attempts to Reverse-Engineer the Formula}
\label{sec:reverse}
%=============================================================================

\subsection{The ratio $\kappa/\beta$}

\begin{equation}
\frac{\kappa_{U(1)}}{\beta_{U(1)}} = \frac{7.269986}{3.796945} = 1.914693\ldots
\end{equation}

This is NOT a recognizable mathematical constant:
\begin{center}
\renewcommand{\arraystretch}{1.2}
\begin{tabular}{lrl}
\toprule
Candidate & Value & Match? \\
\midrule
$35/18$ & 1.94444 & No ($+1.6\%$) \\
$\pi/\varphi$ & 1.94161 & No ($+1.4\%$) \\
$2\ln 2$ & 1.38629 & No \\
$\pi^2/5$ & 1.97392 & No \\
$e/\sqrt{2}$ & 1.92210 & No ($+0.4\%$) \\
\bottomrule
\end{tabular}
\end{center}

The fact that $\kappa/\beta$ is not a recognizable algebraic number
is consistent with it arising from a non-perturbative spectral computation.

\subsection{Testing $\kappa_{U(1)} = \Lambda^3_{\mathrm{full}} \times w \times \varepsilon_{\mathrm{adj}} \times C$}

With $\Lambda^3_{\mathrm{full}} = 885.9375$ (the extended cutoff from
$k = 60$, $a = 9$, $n = 5$, $N = 3$):
\begin{equation}
C_{\mathrm{required}} = \frac{\kappa_{U(1)}}{\Lambda^3_{\mathrm{full}} \cdot w \cdot \varepsilon_{\mathrm{adj}}}
= \frac{7.26999}{885.9375 \times 0.0875 \times 1.000244}
= 0.09376\ldots
\end{equation}

This is tantalizingly close to $3/32 = 0.09375$, but differs by $\sim 10$~ppm.
If $C = 3/32$ exactly, the result would be $\alpha^{-1} = 137.022$
(residual $-104$~ppm), which is insufficient for sub-ppm accuracy.

\subsection{The tree-level formula}

If one naively identifies $\kappa = \beta$ (no renormalization), then:
\begin{equation}
\alpha^{-1}_{\mathrm{tree}} = 6\pi \times \beta_{U(1)}
= 6\pi \times 3.79695 = 71.58\,.
\end{equation}
This is wildly wrong (off by factor $\sim 2$), confirming that
non-perturbative renormalization is essential.

%=============================================================================
\section{Definitive Conclusion}
\label{sec:conclusion}
%=============================================================================

\begin{finding}[Status of the closed-form expression]
\mbox{}
\begin{enumerate}
\item \textbf{The exact electroweak formula is:}
\begin{equation}
\alpha^{-1} = 6\pi\,\kappa_{U(1)}
\end{equation}
This is algebraically exact and derivable from gauge-emergence geometry.

\item \textbf{The $\kappa_{U(1)}$ prediction has two routes:}
\begin{itemize}
\item \textbf{Lattice MC} ($\sim 1\%$): Non-perturbative measurement.
  No closed-form $\beta \to \kappa$ formula.
\item \textbf{Spectral action} (sub-ppm): Numerical evaluation of the
  $a_4$ Seeley--DeWitt coefficient on Toeplitz-truncated
  $\mathbb{C}P^2 \times S^3$.  \textbf{No closed-form algebraic
  expression is written in the DFD corpus.}
\end{itemize}

\item \textbf{The best approximate closed-form is:}
\begin{equation}
\alpha^{-1} \approx \frac{35\pi}{3}\;\beta_{U(1)}\;\frac{d{-}1}{d}
\;\frac{d^2}{d^2{-}1}
= 137.024\ldots
\qquad \text{(85~ppm residual)}
\end{equation}
where $\beta_{U(1)} = 3.796945\ldots$ is the computable 60-term
Chern--Simons sum and $d = 64$.

\item \textbf{The gap:} The full sub-ppm computation
involves a numerical evaluation of the spectral action that has not been
reduced to a single algebraic expression. The paper presents the
\emph{inputs} and the \emph{result} but not the intermediate
computation. This is an \textbf{exposition gap}, not necessarily a
physics gap --- the computation may well be correct but simply not
written out in closed form.
\end{enumerate}
\end{finding}

\paragraph{What would close the gap.}
An explicit evaluation of the $a_4(P)$ coefficient for the connection
Laplacian on $\mathbb{C}P^2 \times S^3$ (7-dimensional, Fubini--Study
$\times$ round metric, with U(1) line bundle twist proportional to the
K\"ahler form), incorporating the Toeplitz truncation at level $d = 64$
and the democratic-to-physics trace conversion
($\varepsilon_{\mathrm{adj}} = 4096/4095$), would provide the missing
closed-form expression. This is a well-posed mathematical computation
in spectral geometry.

%=============================================================================
\section*{File Provenance}
%=============================================================================

This document was produced by exhaustive reading of:
\begin{itemize}
\item \texttt{section\_alpha\_locked.tex} (v108, Section 8C)
\item \texttt{appendix\_K.tex} (v108, Appendix K)
\item \texttt{section\_quantum.tex}, \texttt{section\_alpha\_relations.tex}
\item \texttt{section\_gauge\_coupling.tex}
\item \texttt{appendix\_Z\_complete\_params.tex}
\item \texttt{supplementary/PRL\_SM\_from\_Topology.tex}
\item \texttt{supplementary/DFD\_Standard\_Model\_pasted\_note.tex}
\item ``Ab Initio Derivation of the Fine Structure Constant from DFD'' (PDF, all 26 pages)
\item Prior cross-reference documents:
  \texttt{Alpha\_Inverse\_Formula.tex},
  \texttt{Lattice\_Beta\_Kappa.tex},
  \texttt{CS\_Eigenvalues.tex},
  \texttt{a4\_SeeleyDeWitt.tex},
  \texttt{Kappa\_U1\_Formula.tex},
  \texttt{compute\_alpha.py}
\end{itemize}

\end{document}
